% ============ 02. Marco ============
\section{Marco del ejercicio de validación}

\subsection{Instrumento evaluado}

Se revisó la rúbrica oficial \emph{``Consultoría Estadística · USTA 2026-II ·
Primera entrega (Anteproyecto)''} (PDF de 15 págs. aportado por el equipo). Sus
reglas centrales:

\begin{itemize}
  \item \textbf{8 criterios} en dos bloques: Documento escrito (40\,\%: D1 12\,\%,
        D2 9\,\%, D3 8\,\%, D4 6\,\%, D5 5\,\%) y Sustentación oral (60\,\%:
        S1 28\,\% individual, S2 18\,\%, S3 14\,\%).
  \item \textbf{Criterios bloqueantes: D1 y D3.} Si cualquiera queda en
        \emph{Insuficiente}, el anteproyecto \textbf{se devuelve} aunque la nota
        global pase de 3,0.
  \item \textbf{Referencias inventadas}: una sola referencia inexistente, o que
        no diga lo que se afirma, deja \textbf{D2 en Insuficiente}.
  \item \textbf{Techo de autoría}: nota individual $\leq$ nota de S1 $+$ 1,0.
  \item \textbf{Documento}: 6--10 páginas (sin portada ni referencias), PDF,
        citación \textbf{APA 7}, estructura por secciones ligada a cada criterio.
  \item \textbf{IA}: permitida y \emph{declarada} (media página: en qué se usó,
        cómo se verificó, qué se descartó). No declararla, o declararla en
        falso, deja \textbf{S3 en Insuficiente}.
  \item \textbf{Modalidad}: datos abiertos (este caso) — pertinencia
        argumentada (a quién le importa y por qué, con fuente que documente la
        magnitud en Colombia), procedencia del dato con enlace y licencia,
        demostrar que no es un ejercicio ya resuelto con ese dataset.
\end{itemize}

\subsection{Objeto supuesto del anteproyecto}

Los materiales del equipo documentan una consultoría \emph{ya ejecutada}
(2026-I): la auditoría del repositorio Observatorio MinCiencias (sistema de
reconocimiento de investigadores de MinCiencias, datos abiertos). Se asume que
el anteproyecto 2026-II se construye sobre ese tema en modalidad datos
abiertos; \textbf{esta decisión de objeto debe confirmarla el equipo} (es el
primer punto de la sección de decisiones pendientes).

\subsection{Metodología de la validación}

Tres evaluadores independientes aplicaron la rúbrica sobre los materiales:
(i) el consultor (este documento), (ii) un agente calificador de los criterios
de documento D1--D5 y (iii) un agente calificador de sustentación/formato
S1--S3. Sus informes íntegros:
\texttt{documentacion\_auditoria/data/validacion/validacion\_anteproyecto\_D.md}
y \texttt{...\_S.md}. Los niveles se asignan sobre \emph{lo que hoy existe
redactado} (no sobre el potencial de los materiales).
